\documentclass[10pt,a4paper]{amsart}
\usepackage{amsmath,amssymb,amsthm,graphicx}
% Retained minimal preamble for the Stacks Project chapter more-algebra.tex (More on Algebra).
% The chapter file pinned-original-source.tex carries only \input{preamble} on its first line, so no
% preamble text lives inside the chapter itself.  The Stacks Project uses one single shared preamble for
% every chapter; the snapshot stages it identically as sources/stacks-chow/preamble.tex and
% sources/stacks-etale-cohomology/preamble.tex.  The block below is copied mechanically, line for line,
% from sources/stacks-chow/preamble.tex lines 14-18 (the Xy-pic setup) and 154-217 (the native theorem
% declarations together with the mathematical macros).  Nothing was added, removed or redefined; the
% wrapper class and the amsmath/amssymb/amsthm/graphicx base packages are supplied by the assembler.
% Chapter: more-algebra.tex -- More on Algebra.
% For commutative diagrams we use Xy-pic
\usepackage[all]{xy}

% We use 2cell for 2-commutative diagrams.
\xyoption{2cell}

% Theorem environments.
%
\theoremstyle{plain}
\newtheorem{theorem}[subsection]{Theorem}
\newtheorem{proposition}[subsection]{Proposition}
\newtheorem{lemma}[subsection]{Lemma}

\theoremstyle{definition}
\newtheorem{definition}[subsection]{Definition}
\newtheorem{example}[subsection]{Example}
\newtheorem{exercise}[subsection]{Exercise}
\newtheorem{situation}[subsection]{Situation}

\theoremstyle{remark}
\newtheorem{remark}[subsection]{Remark}
\newtheorem{remarks}[subsection]{Remarks}

\numberwithin{equation}{subsection}

% Macros
%
\def\lim{\mathop{\mathrm{lim}}\nolimits}
\def\colim{\mathop{\mathrm{colim}}\nolimits}
\def\Spec{\mathop{\mathrm{Spec}}}
\def\Hom{\mathop{\mathrm{Hom}}\nolimits}
\def\Ext{\mathop{\mathrm{Ext}}\nolimits}
\def\SheafHom{\mathop{\mathcal{H}\!\mathit{om}}\nolimits}
\def\SheafExt{\mathop{\mathcal{E}\!\mathit{xt}}\nolimits}
\def\Sch{\mathit{Sch}}
\def\Mor{\mathop{\mathrm{Mor}}\nolimits}
\def\Ob{\mathop{\mathrm{Ob}}\nolimits}
\def\Sh{\mathop{\mathit{Sh}}\nolimits}
\def\NL{\mathop{N\!L}\nolimits}
\def\CH{\mathop{\mathrm{CH}}\nolimits}
\def\proetale{{pro\text{-}\acute{e}tale}}
\def\etale{{\acute{e}tale}}
\def\QCoh{\mathit{QCoh}}
\def\Ker{\mathop{\mathrm{Ker}}}
\def\Im{\mathop{\mathrm{Im}}}
\def\Coker{\mathop{\mathrm{Coker}}}
\def\Coim{\mathop{\mathrm{Coim}}}

% Boxtimes
%
\DeclareMathSymbol{\boxtimes}{\mathbin}{AMSa}{"02}

%
% Macros for moduli stacks/spaces
%
\def\QCohstack{\mathcal{QC}\!\mathit{oh}}
\def\Cohstack{\mathcal{C}\!\mathit{oh}}
\def\Spacesstack{\mathcal{S}\!\mathit{paces}}
\def\Quotfunctor{\mathrm{Quot}}
\def\Hilbfunctor{\mathrm{Hilb}}
\def\Curvesstack{\mathcal{C}\!\mathit{urves}}
\def\Polarizedstack{\mathcal{P}\!\mathit{olarized}}
\def\Complexesstack{\mathcal{C}\!\mathit{omplexes}}
% \Pic is the operator that assigns to X its picard group, usage \Pic(X)
% \Picardstack_{X/B} denotes the Picard stack of X over B
% \Picardfunctor_{X/B} denotes the Picard functor of X over B
\def\Pic{\mathop{\mathrm{Pic}}\nolimits}
\def\Picardstack{\mathcal{P}\!\mathit{ic}}
\def\Picardfunctor{\mathrm{Pic}}
\def\Deformationcategory{\mathcal{D}\!\mathit{ef}}
\usepackage{cleveref}
\crefname{originallemma}{original lemma}{original lemmas}
\Crefname{originallemma}{Original lemma}{Original lemmas}
\makeatletter
\newlabel{lemma-completion-regular}{{07NY}{source}{}{sourcepointer-1}{}}
\newlabel{lemma-completion-regular@cref}{{[originallemma][0][]07NY}{[1][0][]source}{}{}{}}
\makeatother
\crefname{originallemma}{original lemma}{original lemmas}
\Crefname{originallemma}{Original lemma}{Original lemmas}
\makeatletter
\newlabel{lemma-completion-dimension}{{07NV}{source}{}{sourcepointer-2}{}}
\newlabel{lemma-completion-dimension@cref}{{[originallemma][0][]07NV}{[1][0][]source}{}{}{}}
\makeatother
\crefname{originallemma}{original lemma}{original lemmas}
\Crefname{originallemma}{Original lemma}{Original lemmas}
\makeatletter
\newlabel{lemma-krasner}{{09EI}{source}{}{sourcepointer-3}{}}
\newlabel{lemma-krasner@cref}{{[originallemma][0][]09EI}{[1][0][]source}{}{}{}}
\makeatother
\title{Every finite separable extension of the completion of a discrete valuation field is obtained by base change from a finite separable extension of the field}
\author{The Stacks Project Authors}
\date{stacks-project:more-algebra.tex:local-snapshot-20261012; GNU FDL 1.2 or later}
\begin{document}
\maketitle
\setcounter{subsection}{1}
\noindent\textit{Editorial note.} All delivered mathematical text is the author's own, copied from the pinned Stacks Project chapter file more-algebra.tex; no statement or proof step was rewritten, reordered or removed. The counted statement is the labeled lemma environment at source lines 33813-33820 (canonical tag 09EJ) and the counted proof is the proof environment at source lines 33822-33852, which immediately follows it and is its own proof.\par\medskip
\noindent\textit{Editorial note.} The author states at proof lines 33825-33827 that the proof works more generally for Noetherian local rings whose completion is a local domain of dimension one; that remark is part of the counted proof and is delivered as written.\par\medskip
\noindent\textit{Editorial note.} The parenthetical remark on the theorem of the primitive element at proof line 33831 is the author's own and is retained inside the counted span.\par\medskip
\noindent\textit{Editorial note.} The proof refers to three earlier lemmas of this chapter, two about completions of discrete valuation rings and one being Krasner's lemma; all three are carried as navigation stubs with their roles recorded in the reasons.\par\medskip
\noindent\textit{Editorial note.} The chapter file pinned-original-source.tex carries only an input of the shared preamble on its first line, so no preamble text lives inside the chapter. The retained file sources/stacks-more-algebra/retained-preamble.tex is a mechanical, line-for-line copy of the project's single shared preamble as staged in the snapshot (Xy-pic setup plus the native theorem declarations and mathematical macros). Nothing was added, removed or redefined, and the excerpt serves only as preamble, never as a span coordinate.\par\medskip
\noindent\textit{Editorial note.} The source URL is the source-level Stacks Project URL, as required when more than one card is produced from one source; the per-claim tag is carried by claim\_key and by this note. The counted statement and proof cite no bibliography entry, so this card needs no bibliography.\par\medskip
\noindent\textit{Allowed standard prerequisites (editorial).}\begin{itemize}
\item The completion of a discrete valuation ring is again a discrete valuation ring, in particular a domain. used at proof lines 33823-33825. Makes the completed ring a complete local domain of dimension one, which is the setting the argument is carried out in.
\item Krasner's lemma: a polynomial with a simple root survives a sufficiently small perturbation of its coefficients, the new root being close to the old one in the maximal-adic topology. invoked at proof lines 33841-33844. Produces a root beta in B of the approximating polynomial P\_1 and allows beta to be chosen so that theta - beta lies in pi B.
\item Nakayama's lemma applied to the finite module B over a complete local domain, and the equality of degrees making a surjective map between the two quotients an isomorphism. used at proof lines 33845-33849. Shows that the comparison map from A-completion[t] over P\_1 to B is surjective and hence an isomorphism, which identifies the base change of L with M.
\end{itemize}\medskip
\ifdefined\hypertarget\hypertarget{sourcepointer-1}{}\fi
\noindent\textit{Source reference.} Original lemma 07NY: sources/stacks-more-algebra/pinned-original-source.tex, source lines 10668--10668. The reference is one of the two applied at proof lines 33823-33824 to know that the completion of A is again a discrete valuation ring; the construction of the extension is inside the counted span.\par\medskip
\ifdefined\hypertarget\hypertarget{sourcepointer-2}{}\fi
\noindent\textit{Source reference.} Original lemma 07NV: sources/stacks-more-algebra/pinned-original-source.tex, source lines 10624--10624. The reference is the second of the two applied at proof lines 33823-33824 for the discreteness of the completion; everything done with that fact is inside the counted span.\par\medskip
\ifdefined\hypertarget\hypertarget{sourcepointer-3}{}\fi
\noindent\textit{Source reference.} Original lemma 09EI: sources/stacks-more-algebra/pinned-original-source.tex, source lines 33748--33748. The reference is applied at proof lines 33841-33844 to obtain the root beta of the approximating polynomial P\_1 in B with theta - beta in pi B; the construction of P\_1 and the Nakayama comparison that consume it are inside the counted span.\par\medskip
\begin{lemma}
\label{lemma-approximate-separable-extension}
Let $A$ be a discrete valuation ring with field of fractions $K$.
Let $A^\wedge$ be the completion of $A$ with fraction field $K^\wedge$.
If $M/K^\wedge$ is a finite separable extension, then
there exists a finite separable extension $L/K$
such that $M = K^\wedge \otimes_K L$.
\end{lemma}

\begin{proof}
Note that $A^\wedge$ is a discrete valuation ring too (by
Lemmas \ref{lemma-completion-regular} and \ref{lemma-completion-dimension}).
In particular $A^\wedge$ is a domain. The proof will work more generally
for Noetherian local rings $A$ such that $A^\wedge$ is a local domain
of dimension $1$.

\medskip\noindent
Let $\theta \in M$ be an element that generates $M$ over $K^\wedge$.
(Theorem of the primitive element.)
Let $P(t) \in K^\wedge[t]$ be the minimal polynomial of $\theta$ over
$K^\wedge$. Let $\pi \in \mathfrak m_A$ be a nonzero element.
After replacing $\theta$ by $\pi^n\theta$ we may assume that
the coefficients of $P(t)$ are in $A^\wedge$. Let
$B = A^\wedge[\theta] = A^\wedge[t]/(P(t))$. Note that $B$ is
a complete local domain of dimension $1$ because it is finite over $A$ and
contained in $M$. Since $M$ is separable over $K$ the element $\theta$
is not a root of the derivative of $P$. For any integer $n$ we can find
a monic polynomial $P_1 \in A[t]$ such that $P - P_1$ has coefficients in
$\pi^nA^\wedge[t]$. By Krasner's lemma (Lemma \ref{lemma-krasner}) we see that
$P_1$ has a root $\beta$ in $B$ for $n$ sufficiently large.
Moreover, we may assume (if $n$ is chosen large enough)
that $\theta - \beta \in \pi B$. Consider the map
$\Phi : A^\wedge[t]/(P_1) \to B$ of $A^\wedge$-algebras
which maps $t$ to $\beta$. Since
$B = \pi B + \sum_{i < \deg(P)} A^\wedge \theta^i$, the map $\Phi$
is surjective by Nakayama's lemma. As $\deg(P_1) = \deg(P)$ it
follows that $\Phi$ is an isomorphism. We conclude that the ring
extension $L = K[t]/(P_1(t))$ satisfies $K^\wedge \otimes_K L \cong M$.
This implies that $L$ is a field and the proof is complete.
\end{proof}

\end{document}
